\scanpage{151}
\mspage{25}
% archive sheet number: 150

\[
\left|\begin{aligned}
&\text{Les }M^\pi=\underline{\mathrm{Filt}}^\pi\text{ sont stables par sommes}\\
&\text{amalgamées simples, et par }\underrightarrow{\lim}\text{ filtrantes de cardinal}\\
&\leq\pi\text{ (pas néc. } \unsure{\text{grandes devant}|\text{grands devant}}\ \pi_0),
\end{aligned}\right.
\tag*{(3.13)}
\]
compte tenu que $r\leq s$ donc $M^r\subset M^s$, donc
$A,B,X$ dans $X_D=\deleted{\unsure{A}}B\amalg_A X$ sont dans $M^s$. Les sommes
amalgamées partielles \uline{finies} des $X_D$ sont
donc aussi dans $M^s$, par récurrence sur le
cardinal de l'ens d'indices, \deleted{\unsure{d'où}} la somme
amalgamée infinie apparaît ainsi comme
$\underrightarrow{\lim}$ filtrante d'objets de $M^s$, l'ens d'indices
étant l'ens. $\mathfrak P_f^*\deleted{\unsure{\subset}}\Delta(C_0,f)$ des parties finies
non vides de $\Delta(C_0,f)$. Cet ens. est de
cardinal majoré par \deleted{$\operatorname{Card}\Delta(C_0,f)\mathbin{\cdot}\unsure{\aleph}$} $s'$, \deleted{dès
que} si $s'$ est un cardinal \uline{infini} tel que
$\operatorname{card}\Delta(C_0,f)\leq s'$. On peut donc prendre $s'=s^{(r^{\pi_0})}$.
On a donc prouvé, compte tenu de (3.13)
appliqué à $M^{s'}$, le

\uline{Lemme} 3.3. On suppose (3.10) i.e. $C_0\subset\mathrm{Fl}^rM$,
$f\in\mathrm{Fl}^sM$, avec \deleted{$r\unsure{\leq}s$} \inserted{$s\geq r\geq\operatorname{Sup}(\pi_1,C_1)$},
et de plus que $C_0$ satisfait la condition de non pléthoricité (3.11.)
(ce qu'on peut toujours supposer, quitte à réduire
$C_0$ sans changer $\widetilde C_0=\Phi$ \deleted{\unsure{ni}} et $C_{0*}=\Psi$), et
que la filtration cardinale satisfait la condition
de stabilité (\unsure{anodine} !) (3.13.). Posons
\[
s'=s^{(r^{\pi_0})}.
\]
Alors $X_1(C_0,f)$ est dans $M^{s'}$. Plus précisément et
plus généralement, pour $C_0$ et $r$ fixés (avec $C_0\subset\mathrm{Fl}^rM$
et satisfaisant la condition de non pléthoricité
(3.11) relativement à $r$), pour \deleted{tout} cardinal
